\subsection{与 Coxeter 系统的关系}

定理 2. \((\mathrm{W},\mathrm{S})\) 是一个 Coxeter 系统。此外，对于 \(s\in \mathrm{S}\) 和 \(w\in \mathrm{W}\)，关系 \(\mathrm{C}(sw) = \mathrm{C}(s).C(w)\) 与 \(l_{\mathbb{S}}(sw) > l_{\mathbb{S}}(w)\) 等价。

对于任意 \(s \in S\)，令 \(P_s\) 为元素 \(w \in W\) 的集合，满足

\[
\mathrm{C} (s). \mathrm{C} (w) = \mathrm{C} (s w).
\]

我们将验证 \(P_{s}\) 满足第 1 节第 7 号的条件 (\(A'\))、(\(B'\)) 和 (C)；然后定理的两个断言将由第 1 节第 7 号的命题 6 得出。

条件 \((\mathbf{A}^{\prime})\) 是显然的。

我们验证 (B')。若 \(P_{s}\) 和 \(sP_{s}\) 有公共元素 w，则有 \(w \in P_{s}\) 和 \(sw \in P_{s}\)，从而

\[
\mathrm{C} (s). \mathrm{C} (w) = \mathrm{C} (s w), \quad \mathrm{C} (s). \mathrm{C} (s w) = \mathrm{C} (w).
\]

于是 \(\mathrm{C}(s).\mathrm{C}(s).\mathrm{C}(w) = \mathrm{C}(w)\)，并且由公式 (5) 可知 \(\mathrm{C}(w) = \mathrm{C}(sw)\)，这与定理 1 矛盾。

我们验证 (C)。设 \(s, s' \in S\)，且 \(w, w' \in W\) 满足 \(w' = ws'\)。假设 \(w \in P_s\) 而 \(w' \notin P_s\)，则

\[
\mathrm{C} (s w) = \mathrm{C} (s). \mathrm{C} (w)\tag{9}
\]

\[
\mathrm{C} (w ^ {\prime}) \subset \mathrm{C} (s). \mathrm{C} (w ^ {\prime})\tag{10}
\]

由 (3) 得。

由 (9) 以及关系 \(w = w's'\)，得到

\[
\mathrm{C} (s) w ^ {\prime} s ^ {\prime} \mathrm{B} = \mathrm{C} (s w).\tag{11}
\]

由公式 \((2')\)，有 \(C(w')\)。\(C(s') \subset C(w') \cup C(w's')\)，立即推出

\[
\mathrm{C} (w ^ {\prime}) s ^ {\prime} \mathrm{B} \subset \mathrm{C} (w s ^ {\prime}) \cup \mathrm{C} (w).\tag{12}
\]

由于 \(C(w')\) 是左陪集 gB 的并，且

\[
\mathrm{C} (s). \mathrm{C} (w ^ {\prime}) = \mathrm{C} (s) w ^ {\prime} \mathrm{B},
\]

公式 (10) 表明 \(C(s)w'\) 与 \(C(w')\) 相交，从而当然 \(C(s)w's'B\) 与 \(C(w')s'B\) 相交。由公式 (11) 和 (12) 可知，双陪集 \(C(sw)\) 等于双陪集 \(C(ws')\) 和 \(C(w)\) 中的一个；由于 \(sw \neq w\)，由定理 1 可得 \(sw = ws'\)。

推论 1. 设 \(w_{1}, \ldots, w_{q} \in \mathbf{W}\)，并令 \(w = w_{1} \ldots w_{q}\)。若

\[
l _ {S} (w) = l _ {S} \left(w _ {1}\right) + \dots + l _ {S} \left(w _ {q}\right),
\]

则

\[
\mathrm{C} (w) = \mathrm{C} (w _ {1}) \dots \mathrm{C} (w _ {q}).
\]

取各 \(w_{i}\) 的约化分解后，问题归结为约化分解的情形

\[
w = s _ {1} \dots s _ {q}, \quad \text { with } s _ {i} \in S.
\]

若 \(u = s_2 \ldots s_q\)，则 \(w = s_1u\)，且 \(l_{\mathbb{S}}(s_1u) > l_{\mathbb{S}}(u)\)，所以由定理有 \(\mathrm{C}(w) = \mathrm{C}(s_1). \mathrm{C}(u)\)。对 \(q\) 作归纳，即得所需公式。

推论 2. 设 \(w \in W\)，令 \(T_w\) 为 \(W\) 的子集，按照第 1 节第 4 号的引理 2 与 \(w\) 相关联。若 \(t \in T_w\)，则

\[
\mathrm{C} (t) \subset \mathrm{C} (w). \mathrm{C} (w ^ {- 1}).
\]

若 \(t \in \mathrm{T}_w\)，则按定义，存在 \(w', w'' \in \mathrm{W}\) 和 \(s \in \mathrm{S}\)，使得

\[
w = w ^ {\prime} s w ^ {\prime \prime}, l _ {S} (w) = l _ {S} (w ^ {\prime}) + l _ {S} (w ^ {\prime \prime}) + 1 \quad \text {and} \quad t = w ^ {\prime} s w ^ {\prime - 1}.
\]

由推论 1，

\[
\mathrm{C} (w). \mathrm{C} (w ^ {- 1}) = \mathrm{C} (w ^ {\prime}). \mathrm{C} (s). \mathrm{C} (w ^ {\prime \prime}). \mathrm{C} (w ^ {\prime \prime - 1}). \mathrm{C} (s). \mathrm{C} (w ^ {\prime - 1}).
\]

因此，

\[
\mathrm{C} (w). \mathrm{C} (w ^ {- 1}) \supset \mathrm{C} (w ^ {\prime}). \mathrm{C} (s). \mathrm{C} (s). \mathrm{C} (w ^ {\prime - 1}).
\]

由 (4)，\(\mathrm{C}(s) \subset \mathrm{C}(s). \mathrm{C}(s)\)。因此，

\[
\mathrm{C} (w). \mathrm{C} (w ^ {- 1}) \supset \mathrm{C} (w ^ {\prime}). \mathrm{C} (s). \mathrm{C} (w ^ {\prime - 1}) \supset \mathrm{C} (t).
\]

推论 3. 设 \(w \in W\)，令 \(H_w\) 为 \(G\) 的子群，由 \(C(w).C(w^{-1})\) 生成。则：

\begin{enumerate}
\def\labelenumi{\alph{enumi})}
\tightlist
\item
  对任意约化分解 \((s_1,\ldots ,s_q)\) 的 \(w\)，
\end{enumerate}

\[
\mathrm{C} (s _ {j}) \subset \mathrm{H} _ {w} \quad \text { for } 1 \leqslant j \leqslant q.
\]

\begin{enumerate}
\def\labelenumi{\alph{enumi})}
\setcounter{enumi}{1}
\tightlist
\item
  群 \(\mathrm{H}_w\) 包含 \(\mathrm{C}(w)\)，且由 \(\mathrm{C}(w)\) 生成。
\end{enumerate}

我们对 \(a)\) 作关于 \(j\) 的归纳证明。假设 \(C(s_k)\) 包含于 \(H_w\)，其中 \(k < j\)。令

\[
t = (s _ {1} \dots s _ {j - 1}) s _ {j} (s _ {1} \dots s _ {j - 1}) ^ {- 1}.
\]

元素 \(t\) 属于定义于第 1 节第 4 号的引理 2 中的子集 \(\mathrm{T}_w\)，它是 \(\mathbf{W}\) 的子集。由推论 2，\(\mathrm{C}(t) \subset \mathrm{H}_w\)，从而 \(\mathrm{C}(s_j) \subset \mathrm{H}_w\)。

由于 \(\mathrm{C}(w) = \mathrm{C}(s_1) \ldots \mathrm{C}(s_q)\)，参见推论 1，故 \(\mathrm{C}(w) \subset \mathrm{H}_w\)，于是得到 \(b)\)。

例。将定理 2 应用于第 2 号所述的 Tits 系统，可知对称群 \(S_{n}\) 以及相邻元素的换位集合构成一个 Coxeter 群。

